% Please do not change %
\documentclass[11pt]{article} % font size
\usepackage[margin=1in]{geometry} % margins
\usepackage{amsmath,amsthm,amssymb} % for the  common "math symbols"
\usepackage{algorithm}
\usepackage{algpseudocode}
\usepackage[shortlabels]{enumitem} % for enumeration
\usepackage{booktabs} % if you need tables
\usepackage{listings}
\usepackage{color}
\DeclareMathOperator*{\argmax}{argmax}
\DeclareMathOperator*{\argmin}{argmin}

% \theoremstyle{definition}
\newtheorem{theorem}{Theorem}[section]
\newtheorem{corollary}[theorem]{Corollary}
\newtheorem{lemma}[theorem]{Lemma}
\newtheorem{proposition}[theorem]{Proposition}
\newtheorem{conjecture}[theorem]{Conjecture}
\newtheorem{definition}[theorem]{Definition}
\newtheorem{fact}[theorem]{Fact}

\lstset{
  frame=none,
  xleftmargin=2pt,
  stepnumber=1,
  numbers=left,
  numbersep=5pt,
  numberstyle=\ttfamily\tiny\color[gray]{0.3},
  belowcaptionskip=\bigskipamount,
  captionpos=b,
  escapeinside={*'}{'*},
  language=haskell,
  tabsize=2,
  emphstyle={\bf},
  commentstyle=\it,
  stringstyle=\mdseries\rmfamily,
  showspaces=false,
  keywordstyle=\bfseries\rmfamily,
  columns=flexible,
  basicstyle=\small\sffamily,
  showstringspaces=false,
  morecomment=[l]\%,
}
% Feel free to include additional packages (if needed), but make sure it does not override the margin/font requirements.
%
\setlength{\parindent}{0pt} % no indent for new paragraphs.
\setlength{\parskip}{5pt} % distance between paragraphs, which will be used when you have a blank line in your LaTeX code.


%%%%%%%%%%%%%%%%%%%%%%%%%%%%%%%%%%%%%%%%%%%%%%%%%%%%%%%%%%%%%%%%%%%%%%%%%%%%%%
\title{Categories, Proofs, and Processes: \\Assignment 1}
%
\author{Wira Azmoon Ahmad}
\date{22 Jan 2024}
%
\begin{document}
%
\maketitle

%You may use the usual \LaTeX\ environments to structure your answers:
%\begin{align}
%        \label{basegnn}
%        \mathbf{h}_u^{(t+1)} = \sigma \bigg( \mathbf{W}_\text{self}^{(t)} \mathbf{h}_u^{(t)}  + \mathbf{W}_\text{neigh}^{(t)} \sum_{v \in N(u)} \mathbf{h}_v^{(t)} + \mathbf{b}^{(t)} \bigg): 
%\end{align}


%\begin{table}[h]
%\caption{A simple table.}
%\centering
%\begin{tabular}[t]{lrrrr}
%\toprule
%Graphs & Can & Be  & Fun & Too \\ 
%\midrule 
%$G_1$ & $1$  & $2$ & $3$ & $4$ \\ 
%$G_2$ & $-1$  & $-2$ & $-3$ & $-4$ \\ 
%$G_3$ & $0.1$  & $0.2$ & $0.3$ & $0.4$ \\ 
%\bottomrule
%\end{tabular}
%\label{tab}
%\end{table}

\section*{Question 1}
\addtocounter{section}{1}
For any computable functions $f,g$ of $L$, their composition $g\circ f$ is also computable and hence in $L$.
The identity function can be defined for each data type, so similar to the set-theoretic view of functions,
the associativity and identity axioms hold for the composition of computable functions.

\section*{Question 2}
The collections $\text{Ob}(\mathcal{C}),\text{Ar}(\mathcal{C})$ are already defined,
along with the domain and codomain as the sets in the binary relations.

\begin{proposition}
    \textbf{Rel} composition satisfies the associativity axiom.
\end{proposition}
\begin{proof}
    Let $Q:W\rightarrow X, R:X\rightarrow Y, S: Y\rightarrow Z$ be arbitrary relations, 
    so $Q\subseteq W\times X, R\subseteq X \times Y, S\subseteq Y \times Z$.
    Then 
    \begin{equation}
    \begin{aligned}
        S\circ (R\circ Q)(w, z) &\iff \exists y: (R\circ Q)(w,y) \wedge S (y, z)\\
        &\iff \exists y: (\exists x: Q(w,x) \wedge R(x,y)) \wedge S(y,z)\\
        &\iff \exists x,y: Q(w,x) \wedge R(x,y) \wedge S(y,z)\\
        &\iff \exists x,y: Q(w,x) \wedge (R(x,y) \wedge S(y,z))\\
        &\iff \exists x: Q(w,x) \wedge (\exists y: R(x,y) \wedge S(y,z))\\
        &\iff \exists x: Q(w,x) \wedge (S\circ R(x, z))\\
        &\iff (S\circ R) \circ Q(w,z)
    \end{aligned}
    \end{equation}
\end{proof}
\begin{proposition}
    \textbf{Rel} composition satisfies the identity axiom.
\end{proposition}
\begin{proof}
    Let $R: X\rightarrow Y$ be an arbitrary relation. By definition of the identity relation, for set $Z$,
    \[\text{id}_Z = \{(a,a)|a\in Z\}\]
    so
    \[R\circ\text{id}_X=R=\text{id}_Y\circ R\]
\end{proof}

\section*{Question 3}
\begin{proposition}
    In \textbf{Set}, an arrow is epic if and only if it is surjective.
\end{proposition}
\begin{proof}
    Suppose $f:X\rightarrow Y$ is surjective, and that $g\circ f=h\circ f$, where $g,h:Y\rightarrow Z$.
    By surjectivity,
    \[\forall y\in Y, \exists x\in X: f(x)=y\]
    Applying either $g$ or $h$, we have
    \[\forall y\in Y, \exists x\in X: g\circ f(x)=g(y)\]
    \[\forall y\in Y, \exists x\in X: h\circ f(x)=h(y)\]
    Since $g\circ f=h\circ f$,
    \[\forall y\in Y, \exists x\in X: g\circ f(x)=g(y)=h\circ f(x)=h(y)\]
    \[\forall y\in Y: g(y)=h(y)\]
    so $g=h$.
    We have 
    \[\forall g,h: g\circ f = h\circ f \Rightarrow g = h\]
    so $f$ is epic.

    Suppose $f$ is not surjective. Then
    \[\exists y\in Y,\forall x\in X: f(x)\neq y\]
    Let $y'\in Y$ be such an element.
    Define $Z=\{0,1\}$, where $g=0$, and $\forall y\in Y-\{y'\},h(y)=0$, with $h(y')=1$.
    Clearly, $g\neq h$ since $g(y')\neq h(y')$.
    Then we have
    \[\forall x\in X: g\circ f(x) = h\circ f(x)=0\wedge f(x)\neq y'\]
    but $g\neq h$.
    Thus, $f$ is not epic.

\end{proof}

\section*{Question 4}
\section*{Question 5}
\section*{Question 6}
\section*{Question 7}




\end{document}
